\chapter{Những câu hỏi khiến thầy muốn thở dài}
\markboth{Những câu hỏi khiến thầy muốn thở dài}{Những câu hỏi khiến thầy muốn thở dài}

\section{Từ câu 901 đến câu 920}
\begin{enumerate}[leftmargin=*]
\item \hsqa{Thầy ơi hôm nay có học không?}{Không, thầy lên lớp để ngắm em thôi.}
\item \hsqa{Thầy ơi mai kiểm tra hả?}{Không, thầy nhắc chơi cho em vui tim.}
\item \hsqa{Bài này có trong vở hả thầy?}{Không, nó tự sinh ra giữa không trung.}
\item \hsqa{Thầy ơi nộp muộn được không?}{Được, nếu thời gian chịu sửa luật riêng cho em.}
\item \hsqa{Thầy ơi học bài cũ hả?}{Ừ, bài mới chưa đủ khổ sao mà quên bài cũ.}
\item \hsqa{Em không đem sách có sao không?}{Không sao, chỉ hơi thiếu dụng cụ đi học thôi.}
\item \hsqa{Thầy ơi bài này chắc dễ mà ha?}{Dễ nếu em chịu học từ trước đó.}
\item \hsqa{Hôm nay mình học tới đâu rồi?}{Tới đoạn thầy thấy em hỏi câu này quá đều.}
\item \hsqa{Thầy ơi có cần ghi không?}{Không cần, nếu em muốn về nhà nhìn đời mông lung.}
\item \hsqa{Em đi vệ sinh được không, ngay lúc thầy phát đề?}{Em chọn thời điểm rất có điện ảnh.}
\item \hsqa{Thầy ơi em không biết mình phải làm gì.}{Bắt đầu bằng nghe hết câu hướng dẫn đi.}
\item \hsqa{Thầy ơi em chưa mang bài, để mai nha.}{Mai của em đông khách thật.}
\item \hsqa{Thầy ơi em ngồi đây được không?}{Được, miễn đó không phải chỗ em gây loạn tốt hơn.}
\item \hsqa{Thầy ơi có cần im lặng không?}{Câu hỏi này tự trả lời luôn đó em.}
\item \hsqa{Em không biết hôm nay trực nhật.}{Lịch trực của em chắc học tàng hình.}
\item \hsqa{Thầy ơi em làm câu này theo cảm giác được không?}{Được, rồi kết quả cũng rất cảm xúc.}
\item \hsqa{Thầy ơi bài kiểm tra này lấy hệ số mấy?}{Hệ số quan tâm của em tăng rất đúng lúc.}
\item \hsqa{Em không có bút, mượn được không?}{Đi học mà quên bút, em đúng là sống phiêu.}
\item \hsqa{Thầy ơi em quên tên bài rồi.}{Tên bài còn quên thì bài nhớ em sao nổi.}
\item \hsqa{Thầy ơi câu này ghi ở đâu?}{Ở chỗ em hay lướt qua nhất đó.}
\end{enumerate}

\section{Từ câu 921 đến câu 940}
\begin{enumerate}[leftmargin=*,resume]
\item \hsqa{Thầy ơi em chép nhầm nên xóa được không?}{Được, miễn em đừng xóa luôn trách nhiệm.}
\item \hsqa{Thầy ơi em không hiểu sao mình bị gọi tên.}{Vì em nổi bật theo hướng không mong muốn.}
\item \hsqa{Bài này dài vậy làm hết hả?}{Không, thầy phát thêm để em ngắm cho vui.}
\item \hsqa{Thầy ơi hôm nay học tiếp nữa hả?}{Ừ, giáo dục chưa có chế độ tua nhanh tới chỗ em thích đâu.}
\item \hsqa{Thầy ơi nộp file ảnh chụp được không?}{Được, nếu ảnh chụp cũng rõ như lời hứa của em.}
\item \hsqa{Em tưởng không cần làm.}{Em tưởng khá nhiều thứ nguy hiểm ghê.}
\item \hsqa{Thầy ơi sao điểm em chưa lên?}{Vì phép màu đang bận.}
\item \hsqa{Em không nghe kịp, thầy nói lại từ đầu đi.}{Thầy nói được, miễn tai em ở lại lần này.}
\item \hsqa{Thầy ơi có cách nào khỏi học phần này không?}{Có, khỏi luôn cơ hội làm đúng phần đó.}
\item \hsqa{Em tưởng nộp online là qua rồi.}{Qua mắt em chứ chưa qua tay thầy.}
\item \hsqa{Sao em bị trừ điểm vậy?}{Vì điểm cũng có tiêu chuẩn chứ không sống cảm tính.}
\item \hsqa{Em học rồi mà sao đề khác.}{Đề có quyền thay áo, em phải hiểu chứ không chỉ nhớ.}
\item \hsqa{Thầy ơi hôm nay mình học nhẹ thôi nha.}{Não em nghe chữ nhẹ là hạnh phúc sớm ghê.}
\item \hsqa{Thầy ơi em quên đem não.}{Mai nhớ mang đủ combo đi học.}
\item \hsqa{Em nghĩ câu này chắc không quan trọng.}{Em hay đánh giá khá tự tin với kiến thức chưa chắc.}
\item \hsqa{Thầy ơi phần này bỏ qua nha.}{Em bỏ, đề không bỏ là khổ.}
\item \hsqa{Em không có thời gian học.}{Lạ, điện thoại em lại có rất nhiều.}
\item \hsqa{Thầy ơi em làm được nhưng không biết giải thích.}{Được nửa đường rồi, nửa còn lại tên là tư duy.}
\item \hsqa{Em chỉ sai một chút.}{Một chút đúng chỗ hiểm là đi nguyên bài.}
\item \hsqa{Thầy ơi có cộng điểm chuyên cần không?}{Có chuyên cần thật rồi tính tiếp em.}
\end{enumerate}

\section{Từ câu 941 đến câu 960}
\begin{enumerate}[leftmargin=*,resume]
\item \hsqa{Em thấy câu hỏi này hơi dư.}{Dư với người chưa biết nó cứu mình chỗ nào.}
\item \hsqa{Em tưởng học sinh được quyền sai.}{Được, nhưng không được quyền lặp sai rồi ngạc nhiên.}
\item \hsqa{Thầy ơi em chưa tải được file.}{Mạng em có duyên né deadline ghê.}
\item \hsqa{Thầy ơi em hỏi lại câu này lần nữa nha.}{Ừ, vì trí nhớ em thương thầy dữ.}
\item \hsqa{Em không thấy thông báo.}{Thông báo chắc thấy em nhưng buồn quá nên đi luôn.}
\item \hsqa{Em nhớ có làm mà.}{Nhớ kiểu mơ hồ này điểm không chấp nhận.}
\item \hsqa{Thầy ơi bài hôm nay có khó không?}{Khó với người chưa chuẩn bị, bình thường với người có não đang bật.}
\item \hsqa{Em tưởng lớp nghỉ sớm.}{Tưởng tượng của em sống phong phú thật.}
\item \hsqa{Thầy ơi sao em học rồi vẫn quên?}{Vì em học qua chứ chưa học vào.}
\item \hsqa{Thầy ơi em tưởng đây là bài cũ.}{Bài cũ mà em còn lạ thì cũng đáng suy nghĩ.}
\item \hsqa{Em viết xấu nên thầy thông cảm.}{Thông cảm có giới hạn, giải mã chữ tượng hình thì không.}
\item \hsqa{Thầy ơi em có thể chọn câu dễ làm thôi không?}{Đời mà cho vậy thì đỡ biết mấy.}
\item \hsqa{Em không biết nộp ở đâu.}{Ở nơi thầy đã nói ba lần và bảng đã ghi một lần.}
\item \hsqa{Em không cố ý bỏ sót.}{Bài của em vẫn bị bỏ sót rất thật.}
\item \hsqa{Thầy ơi em làm bằng suy luận riêng.}{Riêng quá nên đi lạc khỏi đáp án rồi.}
\item \hsqa{Em tưởng câu này không chấm kỹ.}{Giám khảo không làm nghề bỏ qua cho em đâu.}
\item \hsqa{Em không hiểu vì sao thầy lại thất vọng.}{Vì thầy nhìn thấy phần em làm được nhưng cứ chọn làm hời.}
\item \hsqa{Thầy ơi em có thể bắt đầu lại không?}{Có, bài học hay ở chỗ đó luôn còn cửa.}
\item \hsqa{Em chưa sẵn sàng bị kiểm tra.}{Kiểm tra rất hiếm khi hỏi em đã sẵn sàng chưa.}
\item \hsqa{Em chỉ cần thầy tin em thôi.}{Thầy tin nhanh nhất khi thấy em làm thật.}
\end{enumerate}

\section{Từ câu 961 đến câu 980}
\begin{enumerate}[leftmargin=*,resume]
\item \hsqa{Thầy ơi em học rồi nhưng vũ trụ không ủng hộ.}{Vũ trụ bận lắm, em tự lo phần mình đi.}
\item \hsqa{Thầy ơi em thiếu may mắn.}{May mắn của em đang đi lạc vì em chưa chuẩn bị đủ.}
\item \hsqa{Em nghĩ thầy hơi khó với em.}{Khó vừa đủ để em đừng dễ với cái dở của mình.}
\item \hsqa{Thầy ơi sao em học mãi vẫn sai.}{Vì em sửa lỗi chưa sâu, không phải vì thầy trù.}
\item \hsqa{Em chỉ cần thêm chút điểm thôi.}{Chút điểm của em đi kèm khá nhiều chút lười trước đó.}
\item \hsqa{Em thấy mình hơi vô duyên với môn này.}{Vô duyên thì chủ động làm quen chứ đừng đứng xa cười.}
\item \hsqa{Thầy ơi em không làm được vì stress.}{Stress thật cần xử lý, chứ không phải dán lên mọi tờ giấy trắng.}
\item \hsqa{Em nghĩ mình xui từ đầu năm.}{Đầu năm xui hay nếp học em đang xui lâu dài?}
\item \hsqa{Thầy ơi em không dám hỏi vì ngại.}{Ngại một chút đỡ hơn mù lâu dài.}
\item \hsqa{Em tưởng thầy sẽ nhẹ tay.}{Nhẹ tay với lỗi lặp lại là nặng tay với tương lai em đó.}
\item \hsqa{Em học theo cách riêng mà.}{Cách riêng miễn ra kết quả tử tế, không thì nó chỉ riêng thật.}
\item \hsqa{Em không biết vì sao cứ bị thầy nhắc.}{Vì lỗi em chăm chỉ hơn em trong vài việc khác.}
\item \hsqa{Thầy ơi em muốn một câu đỡ buồn.}{Buồn chút để sửa còn hơn vui giả rồi trượt dài.}
\item \hsqa{Em chưa quen cách học này.}{Quen là kết quả của tập, không phải điều kiện mới bắt đầu.}
\item \hsqa{Thầy ơi sao thầy nhớ lỗi em dữ vậy.}{Vì lỗi em ghé lớp đều như điểm danh.}
\item \hsqa{Em nghĩ khác đáp án thì sao?}{Miễn em chứng minh được, đừng chỉ mang cái tôi tới.}
\item \hsqa{Thầy ơi em thấy mình bất lực.}{Bất lực nhất là khi em ngừng thử cách khác.}
\item \hsqa{Em không muốn làm thầy thất vọng.}{Thì bớt để lời hứa chạy trước hành động.}
\item \hsqa{Em muốn được hiểu hơn là bị nhắc.}{Muốn được hiểu thì cũng phải cho người khác thấy em đang cố thật.}
\item \hsqa{Thầy ơi em còn cứu được không?}{Còn, miễn em đừng chỉ hỏi cho cảm động.}
\end{enumerate}

\section{Từ câu 981 đến câu 1000}
\begin{enumerate}[leftmargin=*,resume]
\item \hsqa{Thầy ơi em không biết bắt đầu sửa từ đâu.}{Từ lỗi em ngại nhìn nhất.}
\item \hsqa{Em có nên tin mình làm được không?}{Tin vừa thôi, rồi ngồi làm cho nó đáng tin hơn.}
\item \hsqa{Em thấy mình chậm hơn mọi người.}{Chậm còn kịp sửa, đứng mới đáng sợ.}
\item \hsqa{Em muốn thầy cho lời khuyên cuối.}{Bớt chống chế, tăng nề nếp, nhìn thẳng lỗi mình.}
\item \hsqa{Em có thể làm lại từ bài hôm nay chứ?}{Đó là thứ thầy mong em hỏi nhất.}
\item \hsqa{Em sợ mai lại như cũ.}{Sợ thì tối nay đổi một việc nhỏ trước.}
\item \hsqa{Em chỉ cần ai đó kéo em lên.}{Kéo được, nhưng chân em cũng phải đạp.}
\item \hsqa{Em thấy mình không đủ giỏi.}{Chưa đủ thì bồi, đừng ngồi tự kết án.}
\item \hsqa{Em có hy vọng gì sau chừng này lỗi không?}{Có, vì em vẫn còn biết hỏi và còn nghe.}
\item \hsqa{Thầy ơi điều tệ nhất ở học sinh là gì?}{Không phải dở, mà là dở nhưng mê bênh cái dở.}
\item \hsqa{Điều tốt nhất ở học sinh là gì?}{Biết sai rồi chịu sửa, nhanh hơn người lớn nhiều khi.}
\item \hsqa{Thầy có mệt với tụi em không?}{Có, nhưng mệt không có nghĩa là bỏ.}
\item \hsqa{Thầy có từng bất lực chưa?}{Có chứ, nhưng bất lực không được quyền làm thầy im luôn.}
\item \hsqa{Thầy có từng thấy em hết cứu không?}{Chỉ hết cứu khi em thích ở yên trong lỗi mình.}
\item \hsqa{Thầy nghĩ em nên nhớ câu nào nhất?}{Chậm không sao, đừng đứng.}
\item \hsqa{Em nên giữ gì sau cuốn này?}{Giữ cái tỉnh, bỏ bớt cái cùn.}
\item \hsqa{Em có thể cười vào chính mình không?}{Cười được là tốt, miễn cười xong có sửa.}
\item \hsqa{Thầy ơi cuối cùng học để làm gì?}{Để em bớt sống mù và bớt làm khổ chính mình.}
\item \hsqa{Em còn nhỏ, cần nghĩ sâu vậy không?}{Không nghĩ sớm thì lớn lên nghĩ muộn đau hơn.}
\item \hsqa{Thầy chốt cho em một câu cuối cùng đi.}{Lý do thì ai cũng có, người lớn lên là người bớt lấy lý do làm chỗ trốn.}
\end{enumerate}